% ============================================================
% Chapter 8: Process and Messaging
% ============================================================
\section{进程与消息}

\subsection{进程的注册和获取}

\begin{figure}[H]
\centering
\includegraphics[width=0.85\textwidth]{image48}
\caption{进程注册与获取}
\end{figure}

\subsection{发送信息到远程节点}

向另一节点上的注册进程发送消息的格式：

\begin{lstlisting}
{pong, Pong_Node} ! {ping, self()},
\end{lstlisting}

使用 \lstinline|{registered_name, node_name}| 元组代替单独的 \lstinline|registered_name|。

\begin{figure}[H]
\centering
\includegraphics[width=0.85\textwidth]{image47}
\caption{远程消息发送}
\end{figure}

参考：\href{https://www.erlang.org/doc/getting_started/conc_prog.html#message-passing}{Erlang -- Concurrent Programming}

\subsection{获取远程进程 PID：通过 register}

\begin{figure}[H]
\centering
\includegraphics[width=0.85\textwidth]{image49}
\caption{通过 register 获取远程 PID (1)}
\end{figure}

\begin{figure}[H]
\centering
\includegraphics[width=0.85\textwidth]{image50}
\caption{通过 register 获取远程 PID (2)}
\end{figure}

\begin{figure}[H]
\centering
\includegraphics[width=0.85\textwidth]{image51}
\caption{通过 register 获取远程 PID (3)}
\end{figure}

\subsection{直接用远程 PID 发送消息}

\subsection{io\_request 与 group\_leader}

\begin{figure}[H]
\centering
\includegraphics[width=0.85\textwidth]{image52}
\caption{io\_request 发送到远程}
\end{figure}

\lstinline|io_request| 发送到远程，无法直接执行打印的时候，放到消息队列。

\begin{figure}[H]
\centering
\includegraphics[width=0.85\textwidth]{image53}
\caption{发给 group\_leader 直接执行}
\end{figure}

发给 \lstinline|group_leader| 直接执行，不会放到进程的消息队列中。

\subsection{Observer 观察消息队列}

\begin{figure}[H]
\centering
\includegraphics[width=0.85\textwidth]{image54}
\caption{Observer 消息队列 (1)}
\end{figure}

\begin{figure}[H]
\centering
\includegraphics[width=0.85\textwidth]{image55}
\caption{Observer 消息队列 (2)}
\end{figure}

\subsection{获取进程列表与进程信息}

\begin{itemize}[leftmargin=1.5em]
    \item 获取进程列表
    \item 获取进程信息
\end{itemize}

\subsection{Erlang/OTP 并发编程实践}

\begin{figure}[H]
\centering
\includegraphics[width=0.85\textwidth]{image58}
\caption{并发编程实践}
\end{figure}

修改了一下，只允许一个实例存在：

\begin{figure}[H]
\centering
\includegraphics[width=0.85\textwidth]{image59}
\caption{单实例修改}
\end{figure}
